\documentclass[english, parskip=full]{scrartcl}
\usepackage[utf8]{inputenc}  % Kodierung der Datei
\usepackage[english]{babel}  % Sprache auf Deutsch 
\usepackage{color}
\usepackage{graphicx, gensymb}
\usepackage{amsfonts, cancel, amssymb}
\usepackage{amsmath, amsthm, array, esint}
\usepackage{booktabs, multirow}  % schönere Tabellen
\usepackage{caption}  % Anpassung der Captions
\usepackage{scrlayer-scrpage}    % Manipulation des Seitenstils
\pagestyle{scrheadings}  % Stil für die Seite setzen
\automark{section}  % Abschnittsnamen als Seitenbeschriftung 
\setheadsepline{.5pt}    % Kopzeile durch Linie abtrennen
\usepackage[hidelinks]{hyperref}  % Links und weitere PDF-Features
\usepackage{typearea, tikz, pgffor, verbatim}
\usetikzlibrary{calc, quotes}
\usepackage[cache=false, outputdir=build]{minted}
\usepackage{placeins} % provides \FloatBarrier

\definecolor{bg}{rgb}{0.9,0.9,0.9}
\newcommand\numberthis{\addtocounter{equation}{1}\tag{\theequation}}
\newcommand{\bra}{}
\newcommand{\kom}{}
\newcommand{\brak}{}
\newcommand{\braket}{}
\newcommand{\intx}{}
\newcommand{\intr}{}
\newcommand{\kla}{}
\newcommand{\klam}{}
\newcommand{\klammer}{}
\newcommand{\oper}{}
\newcommand{\tecxt}{}
\newcommand{\Bra}{}
\newcommand{\Ket}{}
\newcommand{\I}{\text{i}}
\renewcommand{\i}{\text{i}}
\newcommand{\e}{\text{e}}
\newcommand*{\Fourier}{\textsc{Fourier}\ }
\newcommand*{\F}{\mathcal{F}}
\newcommand*{\sinc}{\text{sinc\,}}
\newcommand{\conv}{\mathop{\scalebox{3.5}{\raisebox{-0.38ex}{\makebox[0.5\width][c]{$\ast$}}}}\limits}

\newcommand*{\titel}{Variants on Formulae for Frequency Analysis}
\newcommand*{\autor}{Joachim Schwardt}

\hypersetup{pdfauthor={\autor}, pdftitle={\titel}}

%\titlehead{titlehead}
%\subject{subject}
\title{\titel}
%\author{\autor}
\date{\today}

\begin{document}

%\begin{titlepage}
\maketitle  % Titel setzen
%\tableofcontents  % Inhaltsverzeichnis setzen
%\end{titlepage}

\section{Cosine weights}

Consider weights defined by a sum of $K$ $\cos$-functions via
\begin{align}
w_t &\propto \sum_{k=0}^{K-1} a_k\cos \kla\left( {2\pi k \frac{t}{N}} \right),
\end{align}
for a set of coefficients $a_k\in\mathbb{R}$.
Let the proportionality constant be $\frac{1}{C}$, such that the weights are normalized.
Explicitly one finds $C=Na_0$, where $N$ is the signal length.
Then the Fourier coefficients are
\begin{align}
F_j &= \sum_{t=0}^{N-1} c_1\e^{2\pi\i \nu t} w_t \e^{-2\pi\i j\frac{t}{N}} = \\
&= \frac{c_1}{C} \sum_{k=0}^{K-1} \frac{a_k}{2} \sum_{t=0}^{N-1} \e^{2\pi\i\frac{t}{N}(\nu N - j)} \kla\left( {\e^{2\pi\i k \frac{t}{N}} + \e^{-2\pi\i k \frac{t}{N}}} \right) \simeq \\
&\simeq \frac{Nc_1}{2C} \sum_{k=-K+1}^{K-1} a_{|k|}  \intx\int_{0}^{1} \e^{2\pi\i x(\nu N - j - k)} \,\mathrm{d}x = \\
&= \frac{c_1}{4\pi\i a_0} \sum_{k=-K+1}^{K-1} a_{|k|} \frac{\e^{2\pi\i \nu N} - 1}{\nu N - j - k} = \\
&= \frac{c_1}{4\pi\i a_0} \kla\left( {\e^{2\pi\i \nu N} - 1} \right) \sum_{k=-K+1}^{K-1} a_{|k|} \frac{1}{\nu N - j - k}.
\end{align}
Note that the summation must count the $k=0$ term twice! 
The notation is somewhat questionable.

Writing $\nu = \frac{j}{N} + \epsilon$ -- where $j$ is implicitly chosen close enough to the fundamental frequency, such that $\epsilon$ is small --, we obtain the ratio of consecutive coefficients as
\begin{align}
R &:= \frac{F_j}{F_{j+1}} = \frac{\sum_{k=-K+1}^{K-1} a_{|k|} \frac{1}{\epsilon - \frac{k}{N}}}{\sum_{k=-K+1}^{K-1} a_{|k|} \frac{1}{\epsilon - \frac{k+1}{N}}}.
\end{align}
With an FFT one may now numerically evaluate $R$ and solve this equation for $\epsilon$.
Due to the structure of the above, this may become highly unstable if $\epsilon\simeq 0$ or $\epsilon\simeq \frac{1}{N}$.
%We may expand the fraction with the product of all individual terms, i.e.
%\begin{align}
%R &= \frac{\prod_{|m|<K}\sum_{|k|<K} a_{|k|} \frac{\epsilon - m/N}{\epsilon - k/N} }{\prod_{|m|<K}\sum_{|k|<K} a_{|k|} \frac{\epsilon - m/N}{\epsilon - (k+1)/N}}.
%\end{align}
We may expand the fraction with the three critical terms
\begin{align}
R &= \frac{\epsilon - \frac{2}{N}}{\epsilon + \frac{1}{N}} \frac{\epsilon \kla\left( {\epsilon - \frac{1}{N}} \right)  \kla\left( {\epsilon + \frac{1}{N}} \right) \sum_{|k|<K} a_{|k|} \frac{1}{\epsilon - k/N} }{\epsilon \kla\left( {\epsilon - \frac{1}{N}} \right)  \kla\left( {\epsilon - \frac{2}{N}} \right) \sum_{|k|<K} a_{|k|} \frac{1}{\epsilon - (k+1)/N}} = \\
&= \frac{\epsilon - \frac{2}{N}}{\epsilon + \frac{1}{N}} \frac{2a_0 \kla\left( {\epsilon^2 - \frac{1}{N^2}} \right) + 2a_1 \epsilon^2 + \epsilon \kla\left( {\epsilon - \frac{1}{N}} \right)  \kla\left( {\epsilon + \frac{1}{N}} \right)\sum_{1<|k|<K} a_{|k|} \frac{1}{\epsilon - k/N}}{2a_0 \epsilon\kla\left( {\epsilon - \frac{2}{N}} \right) + 2a_1 \kla\left( {\epsilon - \frac{1}{N}} \right)^2 +  \epsilon \kla\left( {\epsilon - \frac{1}{N}} \right)  \kla\left( {\epsilon - \frac{2}{N}} \right) \sum_{1<|k|<K} a_{|k|} \frac{1}{\epsilon - (k+1)/N}}
\end{align}

\subsection{Offset correction}
We may want to correct a median offset of our signal.
This corresponds to $\nu=0$ and therefore gives
\begin{align}
F_j^{(0)} &= \frac{c_0}{Na_0} \sum_{k=0}^{K-1} \frac{a_k}{2} \sum_{t=0}^{N-1} \e^{-2\pi\i j\frac{t}{N}} \kla\left( \e^{2\pi\i k\frac{t}{N}} + \e^{-2\pi\i k\frac{t}{N}} \right) = \\
&= \frac{c_0}{Na_0}  \sum_{k=0}^{K-1} \frac{a_k}{2} \frac{1 - \e^{2\pi\i (k-j)}}{1 - \e^{2\pi\i \frac{k-j}{N}}}
\end{align}
or
\begin{align}
F_j^{(0)} &= \frac{c_0}{Na_0} \sum_{k=0}^{K-1} \frac{a_k}{2} \sum_{t=0}^{N-1} \e^{-2\pi\i j\frac{t}{N}} \kla\left( \e^{2\pi\i k\frac{t}{N}} + \e^{-2\pi\i k\frac{t}{N}} \right) \simeq \\
&\simeq \frac{c_0}{Na_0}  \sum_{|k|<K} \frac{a_{|k|}}{2} \intx\int_{0}^{1} \e^{2\pi\i (k-j)x}\, \mathrm{d}x = \\
&= \frac{c_0}{Na_0}  \sum_{|k|<K} \frac{a_{|k|}}{4\pi\i} \frac{\e^{2\pi\i (k-j)} - 1}{k-j} = 0\quad (\tecxt\text{for } j\neq k).
\end{align}
So let $j=k$.
Then we find the terms
\begin{align}
F_k^{(0)} &= \frac{c_0a_k}{a_0}
\end{align}
Assuming that we are given non-normalized Fourier-coefficients by an algorithm, the proper way to remove this constant offset is by executing the following pseudo-code for all $0<|k|<K$
\begin{align}
|F_k| &\rightarrow |F_k| - |F_0|\cdot\frac{|a_{|k|}|}{2a_0},\\
|F_0| &\rightarrow 0.
\end{align}
The (complex) offset can be computed via $c_0 = \frac{F_0}{Na_0}$

\subsection{Removing peaks}
If we want to extract more than one frequency from a given signal, we need to be able to also remove non-constant peaks.
The complex coefficient for a signal with he frequency $\nu$ and the factor $\epsilon = \nu - \frac{j}{N}$ is given by
\begin{align}
c_1 &= F_j\cdot \frac{4\pi\i}{\e^{2\pi\i\nu N} - 1} \frac{1}{\sum_{k=0}^{K-1} a_k \kla\left( \frac{1}{\epsilon - \frac{k}{N}} + \frac{1}{\epsilon + \frac{k}{N}} \right)} = \\
&= F_j\cdot \frac{4\pi\i}{\e^{2\pi\i N\epsilon} - 1} \frac{\epsilon}{2a_0 + \epsilon\sum_{k=1}^{K-1} a_k \kla\left( \frac{1}{\epsilon - \frac{k}{N}} + \frac{1}{\epsilon + \frac{k}{N}} \right)}.
\end{align}
But in order to remove a peak from the absolute FFT-spectrum, we need to reconstruct the values around the peak-index $j_\tecxt\text{max}$.
This can be done by
\begin{align}
|F_{j_\tecxt\text{max} + j}| &\simeq \left\vert \frac{c_1}{2a_0} \frac{\e^{2\pi\i \nu N} - 1}{2\pi\i} \sum_{k=0}^{K-1} a_k \kla\left( \frac{1}{\epsilon - \frac{k+j}{N}} + \frac{1}{\epsilon + \frac{k-j}{N}} \right) \right\vert = \\
&= \frac{|c_1|\sin( \pi N\epsilon) }{2\pi\epsilon a_0} \left\vert \frac{2a_0\epsilon}{\epsilon - \frac{j}{N}} + a_{|j|} + \sum_{k=-K+1;\, k\neq -j}^{K-1} \frac{\epsilon a_{|k|}}{\epsilon - \frac{k+j}{N}} \right\vert.
%&= \frac{|c_1|\sin( \pi N\epsilon) }{2\pi\epsilon a_0} \left\vert \frac{a_0\epsilon}{\epsilon - \frac{j}{N}} + a_j + \sum_{|k-j|<K;\, k\neq 0} \frac{\epsilon a_{k-j}}{\epsilon - \frac{k}{N}} \right\vert .
\end{align}



%where the $k=0$-term should be counted twice. 
%The magnitude of the coefficients is
%\begin{align}
%|F_j(\nu)| &= \frac{c_1}{2C}\sin\kla\left( {\pi\nu N} \right) \left\vert \sum_{k=-K+1}^{K-1} a_{|k|} \frac{\e^{-\pi\i \frac{k}{N}}}{\sin \klam\left[ {\pi (\nu - \frac{j}{N} + \frac{k}{N})} \right] } \right\vert.
%\end{align}
%Their derivative is
%\begin{align}
%|F_j(\nu)|' &= \pi N\cot (\pi\nu N) F_j +  \frac{c_1}{2C}\sin\kla\left( {\pi\nu N} \right) \left\vert \sum_{k=-K+1}^{K-1} a_{|k|} \frac{\e^{-\pi\i \frac{k}{N}}}{\cos \klam\left[ {\pi (\nu - \frac{j}{N} + \frac{k}{N})} \right] } \right\vert
%\end{align}
%Their derivative is
%\begin{align}
%F_j'(\nu) &= -\pi\i F_j - \frac{2\pi\i N \e^{2\pi\i\nu N}}{1 - \e^{2\pi\i\nu N}}F_j - \frac{c_1}{4C} \dots \sum_{k=-K+1}^{K-1} a_{|k|} \pi \e^{-\pi\i\frac{k}{N}} \frac{\cos \klam\left[ {\pi (\nu - \frac{j}{N} + \frac{k}{N})} \right]}{\sin^2 \klam\left[ {\pi (\nu - \frac{j}{N} + \frac{k}{N})} \right]}.
%\end{align}
%The peak-frequency is the solution to the equation $|F_j'(\nu)| = 0$.

\end{document}